\clearpage\SourceLeaf{87}%
\pdfbookmark[1]{Further Remarks on Determining the Date of Easter}{osterfest1807}
\begin{center}
{\large FURTHER REMARKS ON DETERMINING THE DATE OF EASTER.\par}
\medskip\hrule\medskip
{\small Braunschweigisches Magazin. 1807. September 12.\par}
\medskip\hrule
\end{center}\medskip
The cyclic calculation of Easter has been discussed several times in these pages. It is easy enough in itself, though it rests on principles some readers may be unfamiliar with. It also presupposes auxiliary tables that some readers may lack and, as Superintendent \textsc{Helmuth}’s example has shown, can contain errors even in printed works. Yet it is quite common to want to know the date of Easter for a future year, or for a past year for which no calendar is at hand. Many readers may therefore welcome a special procedure for determining Easter that requires no knowledge of the principles and technical terms of the cyclic method, dispenses with auxiliary tables, and calls for only a couple of very simple operations of arithmetic. I presented this method seven years ago in Mr.~\textsc{von Zach}’s \textit{Monatliche Correspondenz} (the August issue of 1800). Since it will probably have come to the attention of few readers of this magazine, some of them may find it helpful to have the essentials of that procedure explained here again. I shall confine myself here, however, to presenting the rules for determining Easter in the Gregorian calendar for the present century, the immediately %
\SourceLeaf{88}%
following ones and the past ones in a way that even the least experienced reader can understand. I shall also give the rules for the same calculation in the Julian calendar. The extension of the rules to later centuries, and their connection with the original arrangement of our calendar, can be read at the place cited. [P.~75 of this vol.]

To determine Easter for any year of the present century, from 1800 to 1899 inclusive, follow these rules.

$1^{0}$. Divide the year number by 19. Ignore the quotient altogether and record only the remainder. To distinguish it from the remainders that arise in the following operations, we shall call it the \textit{first} remainder. If the division is exact, take 0 as the remainder. The same applies to the other divisions.

$2^{0}$. Multiply this first remainder by 19, add 23 to the product, and divide the sum by 30. Here again record only the remainder, which we call the \textit{second} remainder.

$3^{0}$. Divide the year number by 4. Call the remainder the \textit{third} remainder.

$4^{0}$. Form the sum of four times the year number, six times the second remainder, twice the third, and the number 4. Dividing this sum by 7 gives us the \textit{fourth} remainder.

$5^{0}$. The sum of the second and fourth remainders now shows how many days after the 22nd of March Easter falls. Whenever this sum is at most 9, Easter is still in March, and its date is obtained by adding 22 to that sum. If the sum is greater than 9, Easter falls in April, and its date is obtained by subtracting 9 from the sum.

I shall illustrate these rules with a couple of examples. First, consider the calculation of Easter for the year 1801. Here

\textit{under the first rule}, division by 19 gives the quotient 94, which we can disregard, and the remainder 15, which is therefore the first remainder.

\textit{Under the second rule}, nineteen times the first remainder is 285. Adding 23 gives 308. Dividing this by 30 leaves 8 as the second remainder.

\textit{Under the third rule}, the third remainder is 1.
\SourceLeaf{89}%

\begin{center}\small
\begin{tabular}{llr}
\textit{Under the fourth rule} & we add four times the year number &7204\\
& six times the second remainder &48\\
& twice the third remainder\quad . &2\\
& and the number\quad .\ .\ .\ .\ .\ . &4\\\cline{3-3}
&\hfill giving the sum &7258
\end{tabular}
\end{center}
Dividing this by 7 leaves 6 as the fourth remainder.

\textit{Under the fifth rule}, finally, we have 14 as the sum of the second and fourth remainders. Since this is 5 greater than 9, Easter falls on the 5th of April. — Thus Maundy Thursday of that year, which the events of the day now bring back to mind, fell on the 2nd of April.

For the year 1808, which prompted Superintendent \textsc{Helmuth} to write his paper, the same procedure gives the following.

Under the first rule, the remainder is 3.

Under the second rule, nineteen times the first remainder is 57. Adding 23 gives 80, so the second remainder is 20.

Under the third rule, the third remainder is 0.

Under the fourth rule, we add 7232, 120, 0 and 4. The sum is 7356, so the fourth remainder is 6.

Under the fifth rule, we obtain the sum 26. Subtracting 9 from it gives the 17th of April as the date of Easter.

After these explanations, anyone who knows the four basic operations of arithmetic will be able to carry out the calculation for other years without difficulty.

As given above, these rules apply to the hundred years from 1800 to 1899. For other centuries some changes are needed. In the past century, from 1700 to 1799, the only difference is that in forming the sum under the fourth rule, the number 4 must be replaced by 3. In the seventeenth century, however, or more precisely from 1583 (when Easter was first determined under the new arrangement) to 1699, the fourth rule just mentioned requires that 4 be replaced by 2. In addition, under the second rule, the amount added to 19 times the first remainder must change from 23 to 22. Everything else remains the same. The changes to be made for the centuries immediately ahead also concern these two points alone. Namely, in place of the number 23, which in the %
\SourceLeaf{90}%
first rule applies to the eighteenth and nineteenth centuries, use the number 24 in the twentieth, twenty-first and twenty-second centuries, 25 in the twenty-third, 26 in the twenty-fourth, and 25 again in the twenty-fifth. For the number 4 used in the fourth rule in the nineteenth century, or the 3 used in the eighteenth, take
\begin{center}\small
\begin{tabular}{lr}
in the twentieth and twenty-first\quad .\ .\ .\ . &5\\
in the twenty-second\quad .\ .\ .\ .\ .\ .\ .\ .\ . &6\\
in the twenty-third\quad .\ .\ .\ .\ .\ .\ .\ .\ .\ . &0\\
in the twenty-fourth and twenty-fifth &1
\end{tabular}
\end{center}
The general rules for these changes can be found in the passage cited from the Mon. Corresp. [P.~78 of this vol.]

I have only to add that the epact table in the Gregorian calendar is arranged in such a way that two kinds of exception to the rules above can occur in particular years. The first is that whenever the second remainder is 29 and the fourth is 6, so that by the fifth rule Easter ought to be on the 26th of April, the feast is celebrated one week earlier, on the 19th of April. So far this case has occurred only once, in the year 1609, and it will not recur before 1981. The second exception can occur when the second remainder is 28 and the fourth is 6. By the fifth rule Easter would then fall on the $25^{\mathrm{th}}$ of April, but is actually moved to the 18th of April, and only if the first remainder was not below 11, in which case the exception does not apply and Easter remains on the 25th of April. This second exception has never yet occurred, and will occur for the first time only in the year 1954.

All the foregoing concerns the calculation of Easter in the Gregorian calendar. For the Julian calendar the five rules above require only two changes. In the second rule the number 23 is replaced by 15, and in the fourth rule 4 is replaced by 6. With these changes the rules apply generally to all centuries. I shall illustrate this calculation too with an example. For the year 1808 the first rule, which remains the same, gave us the first remainder 3. Multiplying this by 19 and adding 15 gives 72, so division by 30 leaves %
\SourceLeaf{91}%
12 as the second remainder. The third remainder is 0, as above. Under the fourth rule we now add 7232, 72, 0 and 6. This gives 7310, so the fourth remainder is 2. Adding it to the second remainder gives 14. Subtracting 9 under the fifth rule leaves 5. Easter is therefore on the 5th of April in the Julian style, which is the 17th of April in the Gregorian. So in this year Easter happens to fall on the same day in both calendars, although the Julian calendar gives a date 12 days lower.

\Needspace{9\baselineskip}
It can also be shown that in the 143rd century, that is, from 14200 to 14299, Easter would always fall on the same date in the two calendars if they were to remain in use unchanged for that long, which of course is hardly to be expected. The same date in the two calendars would then always be 15 full weeks apart, so that Julian Easter would come 15 weeks later than Gregorian Easter. This agreement ceases again in the year 14300 and resumes in 43200, when the difference would be 46 weeks.

Finally, I note that the May issue of the Mon. Corresp. for 1802 [p.~80 of this vol.] also gives a purely arithmetical rule for calculating Passover, which can be applied without any further knowledge of how the Jewish calendar is arranged.

\begin{center}\rule{.25\linewidth}{.3pt}\end{center}
